\begin{lemma}[Uniform matching fibers versus hypergeometric masses]
Fix real $0<d\le1$ and $0<\varepsilon<1/16$ with $2\varepsilon<d$.
There is a positive integer $n_0$ such that the following holds uniformly
for all finite graphs $G$, disjoint $U,W$ of size $n\ge n_0$, and
$S\subseteq U$, $T\subseteq W$. Suppose $G[U,W]$ is
$(\varepsilon,d)$-super-regular (parameters embedded into extended
nonnegative reals) and its perfect-matching set is nonempty.
Under \noderef{UniformPerfectMatching}, write $X=|f(S)\cap T|$,
$s=|S|$, $t=|T|$. For every finite set $A$ of natural numbers,
\[
 \Pr(X\in A)\le e^{nL(d,\varepsilon)}\sum_{x\in A}H(n,s,t,x),
\]
with $L$ and $H$ defined by \noderef{MatchingConcentrationCorrection}
and \noderef{HypergeometricMass}. The real nonnegative upper bound is
embedded into extended nonnegative reals in the measure inequality.
\end{lemma}
\begin{proof}
Put $\kappa=\sqrt\varepsilon>0$ and $D=d+\varepsilon+3\kappa>0$.
Embed these finite nonnegative parameters in extended nonnegative reals.
Then $\kappa^2=\varepsilon$, $0<\varepsilon$, $0<d\le1$, and
$2\varepsilon<d$ still hold in that arithmetic.
Take $n_0$ to be at least one and at least both thresholds supplied by
\noderef{MatchingEnumeration} and
\noderef{RegularInducedMatchingUpperBound}, the latter with parameter
$\kappa$. These thresholds depend only on $d,\varepsilon$, not on any
subsets or orders. The super-regular hypothesis supplies regularity by
the definition \noderef{SuperRegularPair}.

Let $M$ be the total matching count. The first cited enumeration statement
gives $M\ge(d-2\varepsilon)^n n!>0$. All these real quantities are finite,
so this is also precisely the real inequality obtained from its
extended-nonnegative-real statement. For any balanced induced rectangle
of order $y$ with left set in $U$ and right set in $W$, the second cited
statement gives at most $D^y y!$ matchings when $y\ge\kappa n$.
For $y<\kappa n$ use the bound $y!$: by
\noderef{PerfectMatching}, every matching is a bijection between its
parts, and there are exactly $y!$ bijections. This elementary count
follows by assigning the first image in $y$ ways and inducting on the
remaining $y-1$ elements; at $y=0$ the unique empty map gives $0!=1$.
Thus the function $K(y)=D^y y!$ for $y\ge\kappa n$, and $K(y)=y!$
otherwise, is nonnegative and bounds all the rectangles required by
\noderef{FourRectangleMatchingFiberBound}.

Consider a feasible $x$, meaning $x\le s,t$ and $t-x\le n-s$.
The four orders in that fiber lemma are
$y_1=x,y_2=s-x,y_3=t-x,y_4=n-s-(t-x)$; they sum to $n$.
Since $\varepsilon<1/16$, $\kappa<1/4$.
Not all four orders can be smaller than $\kappa n$, since their sum
would then be less than $4\kappa n<n$. Let $I$ be the small orders and
$q=\sum_{i\in I}y_i$. Then $|I|\le3$ and $0\le q\le3\kappa n$.
The four rectangle bounds multiply to
\[
 D^{n-q}\prod_i y_i!
 \le D^n\exp(-3\kappa n\log\min(D,1))\prod_i y_i!.
\]
For $D\ge1$, $D^{n-q}\le D^n$ and $\log\min(D,1)=0$.
For $0<D<1$, write $D^{n-q}=D^n\exp(-q\log D)$ and use
$q\le3\kappa n$ and $-\log D>0$. This proves the estimate in both
cases; zero orders contribute $0!=1$ and no density power, as required.

Divide the fiber count from \noderef{FourRectangleMatchingFiberBound}
by $M\ge(d-2\varepsilon)^n n!>0$ and use
\noderef{FourRectangleFactorialCancellation}. This gives
\[
 {M_x\over M}\le
 \left({D\over d-2\varepsilon}\right)^n
 \exp(-3\kappa n\log\min(D,1))
 {\binom t x\binom{n-t}{s-x}\over\binom n s}
 =e^{nL(d,\varepsilon)}H(n,s,t,x).
\]
The equality uses the definitions
\noderef{MatchingConcentrationCorrection} and
\noderef{HypergeometricMass}, and the identity $a^n=e^{n\log a}$
for positive real $a$. Every denominator here is positive.

For an infeasible $x$ no such matching exists: a bijection has at most
$s$ and at most $t$ images in the intersection, while its $t-x$ remaining
preimages of $T$ must fit in $U\setminus S$, of size $n-s$.
The hypergeometric mass is also zero. If $x>s$ this is the explicit
zero clause of its definition; if $x\le s$ but $x>t$, $\binom t x=0$;
otherwise $t-x>n-s$ is equivalent to $s-x>n-t$, and the other binomial
coefficient vanishes. In particular $x>n$ has zero fiber and zero mass.
This proves the pointwise comparison for every natural $x$.

By \noderef{UniformPerfectMatching}, every member of the finite matching
set has mass $1/M$, so the event $X=x$ has probability $M_x/M$.
The events $X=x$ for distinct $x\in A$ are disjoint. Finite additivity
therefore makes $\Pr(X\in A)$ their sum. Sum the pointwise inequality
over $A$ and factor out $e^{nL(d,\varepsilon)}$, obtaining the assertion.
All masses and both finite sums are nonnegative and finite, so embedding
this real calculation into extended nonnegative reals preserves it.
\end{proof}
